\chapter{Betting and Market-Making Games}\label{ch:iv:betting-and-market-making-games}

``Make me a market on the sum of three dice.'' The candidate says 8 at 10. The interviewer buys. ``Again.'' 8 at
10. Buys again. ``Again.'' The fair value is 10.5, and she has sold twice below it, to a counterparty who knows
nothing more than she does and kept buying for exactly one reason. Everything in this chapter is about what the
second ``again'' should have told her: where the fair value is, how wide to quote around it, how much to bet when
the odds are in your favour, and when a trade is information.

\section{Pricing a bet}

The vocabulary of odds and edge is One Quant Book~2's (chapter~29). A bet's fair price is its expected payoff;
its edge is the expected profit per unit staked at the offered price; odds of $b$ to 1 against an event imply a
probability of $1/(b+1)$ at which the bet is fair.

\begin{method}[Pricing any bet in three lines]\label{met:iv:betting-and-market-making-games:price}
\begin{enumerate}
\item List the outcomes with their probabilities and payoffs, net of the stake.
\item Take the expectation; compare it with zero (or with the price asked).
\item Say the edge as a fraction of the stake and the variance, since the size of the bet will depend on both.
\end{enumerate}
\end{method}

\begin{example}[Odds against a six]\label{ex:iv:betting-and-market-making-games:odds}
A bet paying 5 to 1 on a single die showing six wins 5 with probability $\tfrac16$ and loses 1 with probability
$\tfrac56$: expectation 0, fair. At 6 to 1 the edge is $\tfrac16$ of the stake; at 4 to 1 it is $-\tfrac16$.
\end{example}

\section{Sizing}

Knowing a bet is favourable does not say how much to stake. Maximising the expected logarithm of wealth, the
Kelly criterion of One Quant Book~2, chapter~29, gives the fraction $f^\ast = p - q/b$ for a bet that pays $b$ to
1 with probability $p$, and in general the $f$ that solves $\sum_i p_i r_i/(1 + r_i f) = 0$ over outcomes with net
returns $r_i$. Fractions of Kelly trade growth for safety.

\begin{proposition}[Halving under fractional Kelly]\label{prop:iv:betting-and-market-making-games:halving}
In the continuous-time approximation, betting $c$ times the Kelly fraction on a favourable bet, the chance of
ever falling to a fraction $x$ of starting wealth is $x^{2/c - 1}$ (One Quant Book~2, chapter~29). Full Kelly
halves wealth at some point with probability $\tfrac12$; half Kelly with probability $\tfrac18$.
\end{proposition}

\begin{example}[An even-money coin at 60\%]\label{ex:iv:betting-and-market-making-games:kelly}
$f^\ast = 0.6 - 0.4 = 0.2$. Betting 20\% of wealth each time grows wealth at about 2\% a bet in the median, and the
chance of a halving along the way is one half: that is why desks that size by Kelly use a fraction of it.
\end{example}

\section{Making a market on an unknown}

The mechanics are One Quant Book~2's (chapter~30): a two-sided price, a width, and updates as trades arrive.
Interviews test three habits.

\begin{method}[Making a market in an interview]\label{met:iv:betting-and-market-making-games:market}
\begin{enumerate}
\item \emph{Centre} on the expected value, computed or estimated, and say it.
\item \emph{Width} from the uncertainty (a fraction of a standard deviation) and from who may trade with you: wider
if the counterparty may know more.
\item \emph{Update} on information (a card shown, a die revealed) by recomputing the expectation, and on trades by
asking what a counterparty who trades there knows: a trade from someone who knows nothing tells you nothing,
and a trade from someone who knows something tells you what they know.
\item \emph{Track} your position and its risk; say when you are long and by how much.
\end{enumerate}
\end{method}

The hook's candidate failed the first step: 8 at 10 is centred at 9, below the 10.5 of three dice, so an
uninformed counterparty buys at 10 with a half-point edge every time. Moving the market up because of the trade
is the right reflex for the wrong reason; the reason is that the centre was wrong.

\begin{example}[A trade that is information]\label{ex:iv:betting-and-market-making-games:informed}
The interviewer rolls three dice, looks at one and buys whenever that die shows 4 or more. After a buy the fair
value is the mean of that die given it is at least 4, which is 5, plus 7 for the other two: 12. A second buy by
the same player on the same information adds nothing, so the market moves once, not twice.
\end{example}

\section{Informed counterparties and when not to trade}

A market maker who faces some traders who know the value and some who do not earns the width from the second
group and loses to the first, and the balance sets the best width, or says that no width is profitable (Glosten
and Milgrom, 1985; the adverse selection of One Quant Book~1, chapter~1). The chapter's model makes it concrete: the
value is uniform on $[0, 100]$ and the market maker quotes $50 \pm h$; with probability $\alpha$ the arriving trader
knows the value and trades whenever it lies outside the quotes, and otherwise an uninformed trader buys or sells
with probability $\max(0, 1 - h/20)$, less often the wider the market. The expected profit per arrival is
$(1-\alpha)\,h\,(1 - h/20) - \alpha(50 - h)^2/100$ for $h \le 20$.

\begin{omfigure}
\begin{tikzpicture}
\begin{axis}[width=10cm,height=5.4cm,xlabel={\small half-width $h$},ylabel={\small expected profit per arrival},
  xmin=0,xmax=20,ymin=-8,ymax=5,
  tick label style={font=\footnotesize},grid=major,grid style={gray!20},table/col sep=comma,
  legend cell align=left,legend style={font=\scriptsize,at={(0.5,-0.32)},anchor=north,
  /tikz/every even column/.append style={column sep=8pt}},legend columns=3]
\addplot[omDef,thick] table[x=h,y=a05]{figdata/interviews/13-betting-and-market-making-games/width.csv};
\addplot[omProp,thick,dashed] table[x=h,y=a15]{figdata/interviews/13-betting-and-market-making-games/width.csv};
\addplot[omThm,thick,densely dotted] table[x=h,y=a30]{figdata/interviews/13-betting-and-market-making-games/width.csv};
\draw[gray] (axis cs:0,0) -- (axis cs:20,0);
\legend{informed share 5\%,informed share 15\%,informed share 30\%}
\end{axis}
\end{tikzpicture}

\omcaption{A market maker's expected profit per arriving trader against the half-width of the quote, in the
chapter's model (value uniform on $[0, 100]$, uninformed demand falling linearly to zero at $h = 20$). With 5\% or
15\% informed traders the best half-width is about 10.4 or 11.4; with 30\% no width makes money, and the right
answer is not to quote. Exact computation; data: \texttt{fig\_iv\_width.py}.}\label{fig:iv:betting-and-market-making-games:width}
\end{omfigure}

Two related traps complete the family. In a sealed-bid auction for a common value, bidding one's own estimate
loses on average, because winning selects the bidder whose estimate is too high: the winner's curse of One Quant
Book~2, chapter~30, and One Quant Book~4, chapter~29, on common-value auctions and bid shading. And in a game where
one may stop at any time, the value of the option to stop must be computed by \omterm{def:iv:brainteasers-and-logic:backward}{backward induction}, not guessed.

\section{Worked answers}

\begin{example}[A market on a skewed quantity]\label{ex:iv:betting-and-market-making-games:geometric}
\emph{``Make me a market on the number of rolls of a die until the first six.''} The count is geometric with mean 6 and
standard deviation $\sqrt{30} \approx 5.5$, so ``5 at 7'' is centred on the mean. The candidate who stops there misses
what the interviewer is looking for: the distribution is skewed. The median is 4 (the chance of a six within three rolls
is $1 - (5/6)^3 \approx 0.42$, within four $\approx 0.52$). A buyer at 7 loses two times in three (the six comes by the
sixth roll with probability $1 - (5/6)^6 \approx 0.67$) and profits only when it comes after the seventh roll, about 28\%
of the time, but then by a lot. Say so: a counterparty who trades on what usually happens will sell to your bid more
often than buy your offer, and a market centred on the median would be mispriced by two rolls. Then quote, and be ready
to move.
\end{example}

\begin{example}[Locking in a bet that has moved]\label{ex:iv:betting-and-market-making-games:lock}
\emph{``You bet 100 at 3 to 1 on outcome A. The odds are now even money on each side. What do you do if you want no risk,
and what have you made?''} Bet $h$ on B at evens: if A wins you collect 300 and lose $h$; if B wins you lose 100 and
collect $h$. Setting $300 - h = h - 100$ gives $h = 200$ and a locked profit of 100 either way. The profit is the move in
the odds: the first bet implied a 25\% chance and even money implies 50\%, and the hedge converts that repricing into
cash. \emph{Check:} if the odds had not moved (3 to 1 against A, so 1 to 3 on B), the same algebra gives a locked profit
of zero, as it must.
\end{example}

\section{Question bank}

\begin{interviewq}[$\star$ \iqroles{trader} \iqfirm{market maker}]\label{iq:iv:betting-and-market-making-games:1}
You pay 4 to play: a die is rolled and you receive twice its face if it is even and nothing if it is odd. Is the
game fair?
\end{interviewq}

\begin{interviewq}[$\star$ \iqroles{trader} \iqfirm{proprietary firm}]\label{iq:iv:betting-and-market-making-games:2}
A bookmaker offers 3 to 1 against an event. What probability does that imply? If you think the probability is
30\%, what is your edge per unit staked?
\end{interviewq}

\begin{interviewq}[$\star$ \iqroles{trader} \iqfirm{market maker}]\label{iq:iv:betting-and-market-making-games:3}
Two dice are rolled; the contract pays the larger face minus the smaller. Make me a market, and justify the centre and
the width.
\end{interviewq}

\begin{interviewq}[$\star$ \iqroles{trader} \iqfirm{market maker}]\label{iq:iv:betting-and-market-making-games:4}
Make me a market on the number of heads in ten tosses of a fair coin.
\end{interviewq}

\begin{interviewq}[$\star$ \iqroles{trader, researcher} \iqfirm{market maker}]\label{iq:iv:betting-and-market-making-games:5}
What is the fair value of a contract paying the largest of three dice?
\end{interviewq}

\begin{interviewq}[$\star\star$ \iqroles{trader, risk} \iqfirm{proprietary firm}]\label{iq:iv:betting-and-market-making-games:6}
A wager returns three times the stake net with probability 0.35 and loses the stake otherwise. What fraction of your
wealth do you commit each time to maximise long-run growth?
\end{interviewq}

\begin{interviewq}[$\star\star$ \iqroles{trader, researcher} \iqfirm{proprietary firm}]\label{iq:iv:betting-and-market-making-games:7}
Per unit staked, a bet returns $+2$ with probability 0.3, 0 with probability 0.3 and $-1$ with probability 0.4.
What is the Kelly fraction, and what is the growth rate at that fraction?
\end{interviewq}

\begin{interviewq}[$\star\star$ \iqroles{trader, risk} \iqfirm{market maker}]\label{iq:iv:betting-and-market-making-games:8}
You bet repeatedly on a favourable bet. What is the chance your wealth ever halves if you bet the full Kelly
fraction? Half of it? What would you tell a trader who says full Kelly is optimal?
\end{interviewq}

\begin{interviewq}[$\star\star$ \iqroles{trader} \iqfirm{market maker}]\label{iq:iv:betting-and-market-making-games:9}
Make me a market on the number of red cards among the top ten of a shuffled deck. The first three are turned over
and all are red. Where is fair value now?
\end{interviewq}

\begin{interviewq}[$\star\star$ \iqroles{trader} \iqfirm{market maker}]\label{iq:iv:betting-and-market-making-games:10}
Three dice are rolled; the interviewer looks at one of them and buys from you whenever it shows 4 or more. You
quote 10 at 11 and are bought from. Where is fair value? You are bought from again. Where is it now?
\end{interviewq}

\begin{interviewq}[$\star\star\star$ \iqroles{trader, researcher} \iqfirm{market maker}]\label{iq:iv:betting-and-market-making-games:11}
In the chapter's model (value uniform on $[0, 100]$, quote $50 \pm h$, a share $\alpha$ of informed traders,
uninformed demand $\max(0, 1 - h/20)$), what half-width maximises profit when $\alpha = 0.15$, and what is the
profit? What do you do when $\alpha = 0.3$?
\end{interviewq}

\begin{interviewq}[$\star\star\star$ \iqroles{trader, researcher} \iqfirm{proprietary firm}]\label{iq:iv:betting-and-market-making-games:12}
Two firms bid in a sealed-bid auction for an asset worth $V$ to both; each sees $V$ plus independent noise uniform
on $[-10, 10]$, bids its estimate minus $k$, and the higher bid wins and pays its bid. For what $k$ does the
winner break even on average? Why is it not zero?
\end{interviewq}

\begin{interviewq}[$\star\star\star$ \iqroles{trader} \iqfirm{market maker}]\label{iq:iv:betting-and-market-making-games:13}
A deck of 10 red and 10 black cards is turned over one card at a time. Each red pays you 1 and each black costs you
1, and you may stop at any moment. What is the game worth with optimal stopping? Why is it worth more than zero
when the deck is balanced?
\end{interviewq}

\begin{interviewq}[$\star\star\star$ \iqroles{trader, risk} \iqfirm{proprietary firm}]\label{iq:iv:betting-and-market-making-games:14}
With a bankroll of 1\,000 you must bet exactly 100 at even money on each of 100 rounds of a coin that wins with
probability 0.55, stopping only if you are broke. What is the chance of ruin, and what is your expected final
bankroll? How does fixed-size betting compare with betting a fraction of wealth?
\end{interviewq}

\begin{omsources}
\item J.~L.~Kelly, ``A new interpretation of information rate'', \emph{Bell System Technical Journal} 35(4), 1956,
917--926.
\item L.~R.~Glosten and P.~R.~Milgrom, ``Bid, ask and transaction prices in a specialist market with
heterogeneously informed traders'', \emph{Journal of Financial Economics} 14(1), 1985, 71--100.
\item One Quant Book~2, chapters 29--30 (expected value, Kelly, market-making games); One Quant Book~4,
chapter~29 (auctions).
\end{omsources}
